\section{Problem Formalisation}\label{sec:preliminaries}\label{sec:formal}

% Main-text formal part (plan section 9.2: about one page). Wording follows decision D005.
% Full proofs, the fragment assumptions and the IR grammar are in supplement.tex.

Writing message $h_r$ of a history commits version $r$ of the claim graph $\G_r$, a set of
entries of four kinds, each with a version stamp $\vs$:
\emph{claims} $c=(s,p,o,\omega,\tau,\vs)$ with scope $\omega$ and valid interval
$\tau=[\tau^-,\tau^+)$ (start possibly unknown, end possibly open);
\emph{attributions} $(u,c,f,\rho,\pi,\gamma,\delta,\vs)$, a stance of speaker $u$ on $c$ with
communicative function $f$, report mode $\rho$, polarity $\pi$, certainty $\gamma$ and a flag
$\delta$ for derived restatements;
\emph{evidence links} $(c,m,\vs)$ to a message span $m$;
and \emph{supersession events} $(c,\kind,\vend,\vs)$: from version $\vs$ on, $c$ is no longer
currently accepted, because the world changed ($\kind=\change$, optionally with a valid end
$\vend$) or $c$ was wrong ($\kind=\correction$).
The graph is append-only: $\G_{r-1}\subseteq\G_r$, and all entries added at $r$ carry stamp $r$.
Vocabulary and entity identities are assumed fixed between the versions compared; a vocabulary
change or an entity merge invalidates every view.

Queries are built from \emph{claim patterns} $P=(s,p,o,\omega)[u,f]$, whose optional bracket
requires an attribution by $u$ with function $f$ (by default own, positive, not derived): one to
three joined positive patterns, filters against constants, at most one \textsc{not exists} over
one pattern, and a projection, distinct count or latest-state selection (Supplement~S6). Each
pattern carries the built-in negation ``currently accepted'' (or ``not corrected'' for historical
reads). $\cut{\Q}{r}$ sees only entries stamped at most $r$, so $\cut{\Q}{r_0}$ gives the same
result on every $\G_r$, $r\ge r_0$: only reads whose cut-off moves can go stale.
$\ans(\cut{\Q}{r},\G)$ is the answer set; the witnesses $\wit$ of an answer are its minimal sets
of claims and constraint-satisfying attributions over the positive part (why-provenance
\cite{buneman2001why}; a count takes the union over its members, a selection the witnesses of the
selected claim); $\ev$ is the set of evidence spans of the witness claims.

\begin{definition}[Two-level staleness]
At version $r\ge r_0$, a view of $\Q$ with cut-off $r_0$ is \emph{answer-stale} iff
$\ans(\cut{\Q}{r},\G_r)\neq\ans(\cut{\Q}{r_0},\G_{r_0})$, and \emph{citation-stale} iff the
answers are equal but $(\wit,\ev)$ differ.
\end{definition}

\noindent An answer-stale view needs new text; a citation-stale one only has its sentence-level
citations patched, by a program.

\begin{observation}
Invalidating a view only when a claim it used is superseded misses a new match, a write that
supersedes nothing but matches the negated pattern (an acceptance empties ``offered and not yet
accepted''), and a new aggregate member. None of the three supersedes anything: in an
append-only graph every change is an insertion, and new claims and attributions change results
without superseding a used claim.
\end{observation}

\noindent Let $\Dlt=\G_r\setminus\G_{r_0}$. A claim \emph{matches} $P$ if its positions agree
with the constants of $P$ and, if $P$ has a bracket, it carries a satisfying attribution (state,
valid time and filters ignored). An \emph{answer trigger} of $P$ is a claim of $\Dlt$ matching
$P$, an attribution of $\Dlt$ through which its claim matches $P$, or a supersession event of
$\Dlt$ on a claim matching $P$; a \emph{citation trigger} is an evidence link of $\Dlt$ on a claim
matching a positive pattern.

\begin{proposition}[Soundness]\label{prop:soundness}
If no entry of $\Dlt$ is an answer trigger of a pattern of $\Q$, positive or negated, the view is
not answer-stale at $r$ and its witnesses are unchanged; if there is no citation trigger either,
it is not citation-stale.
\end{proposition}

\noindent The rows of a pattern depend only on a claim and the attributions and supersession
events on it, so without a trigger each pattern has the same rows on $\G_{r_0}$ and $\G_r$; the
rest follows by structural induction (Supplement~S3). The argument is routine and extends
Proposition~1 of Martin et al.~\cite{martin2010caching} (one triple, \textsc{and}, \textsc{opt},
\textsc{union}, \textsc{filter}) to negation, distinct counting, latest-state selection, cut-off
reads, claim patterns with attribution constraints, and witnesses and evidence; it stays sound
when decided per write and with each pattern's local filters.

\begin{proposition}[Exactness]\label{prop:exactness}
With local filters, the decision marks a view iff its answers, witnesses or evidence change,
provided $\Q$ has a single pattern, without negation, count, selection, world time, repeated or
scope variables, or undecidable filters or filters on $\mathrm{end}(c)$; and no claim reached by a claim, attribution or
evidence trigger is superseded by $r$, no supersession trigger hits a claim superseded at $r_0$,
and no evidence trigger repeats a cited span (Supplement~S4).
\end{proposition}

\noindent Otherwise the decision over-approximates: a new claim without a join partner, an entry
matching the negated pattern but no answer, a claim with an earlier start under selection by valid
time. Even under the conditions a mark does not tell answer from citation staleness (a second
acceptance changes only the witnesses). Exact instance-level decision is re-evaluation, so
read-time re-query with fingerprints is the reference for every strategy.
